\chapter{The Asia-Pacific Map}\label{ch:nw:the-asia-pacific-map}

India's securities regulator received complaints in 2015 that the National Stock Exchange's \omterm{def:nw:colocation-products-and-how-they-are-sold:colo}{colocation} facility sent its tick-by-tick
data over TCP, member by member, in the order in which members had connected, so that ``the first one to connect to the lowest load
server would get advantage in terms of receiving the data faster than others''. Its 2019 order described the design in detail and
directed the exchange to disgorge 624.89~crore rupees with interest; the appeals tribunal set the disgorgement aside, the regulator
appealed, and on 18 September 2026 the Supreme Court allowed a settlement of about 1\,492~crore rupees that closed the case. No cable
was shorter and no server was faster: the advantage was built into the order in which one computer sent the same message to many.

This chapter maps the Asia-Pacific venues (Tokyo, Hong Kong, Singapore, Sydney, Mumbai) with the method of chapters~10 and~11, and then
studies the case: what sequential dissemination does to the time at which each member sees a tick, what a randomiser and a second server
change, and why multicast ends the problem. Distances in this region are thousands of kilometres, and several venues do not publish the
address of their building, so the site table carries its precision with it.

\section{Tokyo}

JPX runs its co-location in what it calls the JPX Primary Site, where trading participants, information vendors and independent software
vendors install their equipment next to the trading systems of the Tokyo Stock Exchange and the Osaka Exchange. JPX does not publish the
building's address; its network, arrownet, reaches the same trading systems from access points outside the site. The site table therefore places Tokyo at city level, which is an error of kilometres (tens of microseconds)
and irrelevant at the scale of this map: Tokyo is \qty{2880}{\kilo\metre} from Hong Kong, \qty{9.6}{\milli\second} of light in vacuum and
\qty{14.0}{\milli\second} in straight fibre.

\begin{dated}{2026-09}{Where the Asia-Pacific venues match}\label{dat:nw:the-asia-pacific-map:sites}
\begin{itemize}
\item \textbf{Tokyo}: JPX's Primary Site hosts co-location next to the TSE and OSE trading systems; the building is not published (city level).
\item \textbf{Hong Kong}: HKEX's data centre in Tseung Kwan O Industrial Estate (1 Chun Ying Street) hosts its co-location.
\item \textbf{Singapore}: SGX is moving its primary data centre to a new building, which members may use for co-location, together with its new
  trading platform, Iris-ST; its secondary data centre is unchanged; mandatory conformance testing runs from January to April 2027 (city level).
\item \textbf{Sydney}: ASX runs its primary markets and post-trade platforms from the Australian Liquidity Centre, at Gore Hill, since
  February 2012.
\item \textbf{Mumbai}: NSE's co-location is at its Exchange Plaza site in the Bandra Kurla Complex, with over 200 members on racks and about 100
  more through co-location as a service.
\end{itemize}
\end{dated}

\section{Hong Kong and Singapore}

HKEX offers co-location, which it calls hosting services, in its own data centre in the Tseung Kwan O Industrial Estate, on the eastern
side of Hong Kong. SGX is in the middle of a larger move. Its questions and answers on Iris-ST, its new trading platform, say that the
platform will go live in a newly built primary data centre that members can also use for co-location, with the trading engine in a ``2+1''
arrangement (a primary and a hot standby in the primary data centre, a third instance at the secondary site). They also answer a question
this chapter is about: from co-location, members will be able to receive the ITCH market data both over TCP and over multicast.

\begin{table}[ht]
\centering\small
\begin{tabular}{llrrr}
\toprule
From & To & Geodesic & Vacuum & Fibre \\
 & & (km) & (ms) & (ms) \\
\midrule
Tokyo & Hong Kong TKO & 2\,880.3 & 9.61 & 14.05 \\
Hong Kong TKO & Singapore & 2\,577.1 & 8.60 & 12.57 \\
Tokyo & Singapore & 5\,311.0 & 17.72 & 25.90 \\
Singapore & Sydney ALC & 6\,296.4 & 21.00 & 30.71 \\
Tokyo & Sydney ALC & 7\,784.2 & 25.97 & 37.96 \\
Singapore & Mumbai BKC & 3\,903.3 & 13.02 & 19.04 \\
Hong Kong TKO & Mumbai BKC & 4\,317.3 & 14.40 & 21.05 \\
Tokyo & Mumbai BKC & 6\,743.1 & 22.49 & 32.88 \\
\bottomrule
\end{tabular}
\caption{Distances and one-way \omterm{def:nw:the-north-american-map:floor}{latency floors} between the Asia-Pacific sites of \cref{dat:nw:the-asia-pacific-map:sites}. Tokyo and Singapore
are placed at city level, the others at street, neighbourhood or district level: errors of a few kilometres, tens of microseconds at most, against
floors of milliseconds. Data: \texttt{nw\_asia.pair\_rows()}.}\label{tab:nw:the-asia-pacific-map:pairs}
\end{table}

\section{Sydney}

ASX moved its equities trading platform into the Australian Liquidity Centre, a data centre at Gore Hill on Sydney's North Shore, in
February 2012. Its page states the two rules of chapter~9 in one sentence: every customer is given ``identical length cross connects'' to the
trading platform wherever it sits in the building, and the building is carrier-neutral, so that its customers can bring their own networks.
Sydney is the most isolated of the region's centres: \qty{6296}{\kilo\metre} from Singapore (\qty{21.0}{\milli\second} in vacuum) and
\qty{7784}{\kilo\metre} from Tokyo. A trading firm in Sydney competes with other firms in the same building for ASX's own markets, and with
nobody nearby for anything else.

\begin{omfigure}
\begin{tikzpicture}
\begin{axis}[axis equal image,width=11.5cm,xmin=-7200,xmax=5200,ymin=-5000,ymax=4300,xtick distance=2000,xlabel={km east of 120\textdegree{} E},
  ylabel={km north of 5\textdegree{} N},tick label style={font=\footnotesize,/pgf/number format/1000 sep={\,}},label style={font=\small},grid=major,grid style={gray!20},
  table/col sep=comma,unbounded coords=jump,scaled ticks=false]
\addplot[omDef,thick,no markers] table[x=x,y=y]{figdata/networks/12-the-asia-pacific-map/map_links.csv};
\addplot[only marks,mark=*,mark size=2.4pt,omThm,visualization depends on={value \thisrow{anchor}\as\anc},nodes near coords,
  point meta=explicit symbolic,every node near coord/.append style={font=\footnotesize,anchor=\anc,black}]
  table[x=x,y=y,meta=name]{figdata/networks/12-the-asia-pacific-map/map_sites.csv};
\addplot[only marks,mark=none,nodes near coords,point meta=explicit symbolic,
  every node near coord/.append style={font=\scriptsize,fill=white,inner sep=1pt,anchor=center,text=omMeth}]
  table[x=x,y=y,meta=text]{figdata/networks/12-the-asia-pacific-map/map_labels.csv};
\end{axis}
\end{tikzpicture}

\omcaption{The Asia-Pacific venues from computed coordinates, each link labelled with its \omterm{def:nw:the-north-american-map:geodesic}{geodesic distance} and light-time in vacuum. An
equirectangular projection centred on 5\textdegree{}~N, 120\textdegree{}~E stretches shapes badly over these distances: it is for drawing only,
and every number on the map is computed on the ellipsoid. Data: \texttt{fig\_asia.py} on \texttt{firm.geomap}.}\label{fig:nw:the-asia-pacific-map:map}
\end{omfigure}

\section{Mumbai and the colocation case}

\begin{definition}[Unicast dissemination]\label{def:nw:the-asia-pacific-map:unicast}
\emph{Unicast dissemination}\index{unicast dissemination} is the distribution of market data by sending a separate copy of each message to
each recipient, over a connection per recipient (typically TCP), so that one sender transmits the copies one after another.
\end{definition}

NSE's tick-by-tick feed, according to SEBI's order, was disseminated over TCP from June 2010 in the futures and options segment, and
multicast arrived only in April 2014 (November 2014 in the cash market). The order describes the design in two levels. A dissemination
centre sent each batch of ticks to the distribution (POP) servers, one after another, in the order in which the servers had logged in that
day. Each POP server ran three sender processes, called ports, and each port sent to its own list of members, one after another, in the
order in which the members had logged in. A member who logged in first, on a port of a server that had logged in first, was first in both
queues; a member on a lightly loaded server, such as a backup server meant for business continuity, had fewer members ahead of it.

\begin{omfigure}
\begin{tikzpicture}[font=\footnotesize,b/.style={draw,rounded corners=2pt,minimum height=0.55cm,align=center,inner sep=2pt},>=Stealth]
\node[b,fill=omThm!20,minimum height=1.6cm] (c) at (0,0) {dissemination\\ centre};
\foreach \k/\y/\lab in {1/1.8/{server 1\\ (logged in first)},2/0.6/{server 2},3/-0.6/{server 3},4/-1.8/{secondary\\ (lightly loaded)}}{
  \node[b,fill=omDef!15,minimum width=2.2cm] (s\k) at (3.4,\y) {\lab};
  \draw[->] (c.east) -- node[pos=0.55,fill=white,inner sep=1pt,font=\scriptsize]{\k} (s\k.west);}
\foreach \p/\y in {1/2.4,2/1.8,3/1.2}{
  \node[b,fill=gray!10] (p\p) at (6.4,\y) {port \p};
  \draw[->] (s1.east) -- (p\p.west);}
\node[b,fill=white,minimum width=4.3cm,anchor=west] (m1) at (7.4,2.4) {members in login order: 1st, 2nd, \dots, $m$-th};
\draw[->] (p1.east) -- (m1.west);
\node[align=left,font=\scriptsize,anchor=west] at (7.4,1.5) {each port sends the same tick to its members\\ one after another, in the
  order they logged in};
\node[align=left,font=\scriptsize,anchor=west,omMeth] at (5.4,-0.9) {multicast: one copy, replicated by the\\ switches to every
  member at once};
\end{tikzpicture}

\omcaption{The two-level unicast design of NSE's TCP tick-by-tick feed as SEBI's 2019 order describes it, schematically: the centre sends to
its distribution servers in their login order, each server's ports send to their members in the members' login order. Multicast replaces both
queues with one copy that the network replicates.}\label{fig:nw:the-asia-pacific-map:arch}
\end{omfigure}

\begin{proposition}[Sequential unicast]\label{prop:nw:the-asia-pacific-map:seq}
If the centre spends $c_0$ per copy to a server and a port spends $c$ per copy to a member (times a factor $\sigma_s$ for a slower server $s$),
the member at login rank $r$ ($r = 1, 2, \dots$) on a port of the server that ranks $q$ in the centre's order receives a tick after
\[
  t(q, r) = q\,c_0 + r\,c\,\sigma_s + \varepsilon,
\]
with $\varepsilon$ a small random delay. With $m$ members per port, the median member waits about $c\,m/2$ more than the first on the same
port, whatever the network.
\end{proposition}

\begin{proof}
The centre's copies leave one after another, so the $q$-th server receives the tick after $q\,c_0$; the port then sends $r$ copies before
the $r$-th member's. The median rank on a port of $m$ members is about $m/2$.
\end{proof}

\begin{definition}[Dissemination fairness]\label{def:nw:the-asia-pacific-map:fair}
\emph{Dissemination fairness}\index{dissemination fairness} is the property of a market-data distribution that recipients in the same
service receive each message at the same time, up to a spread that does not depend on who they are: measured by the spread of receive
times for one message and by whether the order of recipients persists from one message to the next.
\end{definition}

SEBI's circular of May 2015 put the rule in regulatory form: exchanges must give co-located members ``fair and equal access to facilities
and data feeds'', must ensure that they ``experience similar latency with respect to exchange provided infrastructure'', and must publish
quarterly latency reports. The definition has two parts because the case had two. The spread says how much faster the first recipient is
than the last. The persistence says whether it is always the same one: a random order with a large spread gives everyone a turn at being
first; a fixed order gives one member every race.

\omcode{code/firm/dissem/firm_dissem.py}{47}{64}{Two-level unicast dissemination: the centre's copies to the servers in server order, each port's
copies to its members in login order, and the randomiser that redraws both orders for each tick.}\label{lst:nw:the-asia-pacific-map:unicast}

\begin{omfigure}
\begin{tikzpicture}
\begin{axis}[width=11cm,height=6.2cm,xmin=0.5,xmax=10.5,ymin=0,ymax=130,xlabel={member's login rank on its port},
  ylabel={mean receive delay (\si{\micro\second})},tick label style={font=\footnotesize},label style={font=\small},grid=major,
  grid style={gray!25},legend columns=4,legend style={font=\footnotesize,at={(0.5,-0.25)},anchor=north,draw=none},
  table/col sep=comma,xtick={1,2,3,4,5,6,7,8,9,10}]
\addplot[omDef,thick,mark=*,mark size=1.5pt] table[x=rank,y=s1]{figdata/networks/12-the-asia-pacific-map/rank_curve.csv};
\addplot[omThm,thick,mark=square*,mark size=1.5pt] table[x=rank,y=s2]{figdata/networks/12-the-asia-pacific-map/rank_curve.csv};
\addplot[omMeth,thick,mark=triangle*,mark size=1.8pt] table[x=rank,y=s3]{figdata/networks/12-the-asia-pacific-map/rank_curve.csv};
\addplot[gray,thick,mark=diamond*,mark size=1.8pt] table[x=rank,y=s4]{figdata/networks/12-the-asia-pacific-map/rank_curve.csv};
\legend{server 1st in the centre's order,2nd,3rd,4th}
\end{axis}
\end{tikzpicture}

\omcaption{Simulation, not a measurement of any venue: mean receive delay of a tick by the member's login rank on its port, for the four busy
servers of the model (three ports of ten members each; \qty{5}{\micro\second} per copy at the centre, \qty{10}{\micro\second} per copy at a
port, exponential jitter of mean \qty{1}{\micro\second}; 200 ticks). Each member's delay is fixed by two login orders. Data:
\texttt{fig\_asia.py}, \texttt{nw\_asia.rank\_curve()}.}\label{fig:nw:the-asia-pacific-map:rank}
\end{omfigure}

\Cref{fig:nw:the-asia-pacific-map:rank,tab:nw:the-asia-pacific-map:designs} run the model with parameters chosen for illustration, not taken
from the order. Under plain unicast, the first member receives a tick after \qty{16}{\micro\second} and the median member after
\qty{66}{\micro\second}; the order is the same tick after tick (a rank correlation of 1.00), so the member who was first on one tick beats the
median member on every later one. A randomiser, which the order says NSE had developed in 2011 but not implemented in this feed's dissemination, redraws the order for
every tick: the spread does not change, but the rank correlation falls to 0.07. It still leaves the lightly loaded server: in the model the
first member of a randomised tick is usually one of the secondary server's six, and it beats the median member 88\% of the time. Closing the
secondary to normal traffic brings that to 51\%, a coin toss; multicast shrinks the spread itself from about \qty{108}{\micro\second} to the
jitter.

\begin{table}[ht]
\centering\small
\begin{tabular}{lrrrrr}
\toprule
Design & First & Median & Spread & Rank & First beats \\
 & (\si{\micro\second}) & (\si{\micro\second}) & (\si{\micro\second}) & correlation & median \\
\midrule
Unicast, fixed login order & 16.0 & 65.6 & 107.7 & 1.00 & 100\% \\
Unicast, randomised & 15.1 & 70.7 & 110.2 & 0.07 & 88\% \\
Unicast, randomised, no secondary & 15.2 & 69.0 & 106.9 & 0.01 & 51\% \\
Multicast & 10.0 & 10.7 & 8.4 & -- & 50\% \\
\bottomrule
\end{tabular}
\caption{Simulation of four dissemination designs (the parameters of \cref{fig:nw:the-asia-pacific-map:rank}; 126 members, of whom six on the
secondary server, which the fourth design closes): one tick's first, median and spread of receive times, the rank correlation between
consecutive ticks, and how often the member first on one tick beats the median member of that tick over the next 200 ticks, both reacting in
\qty{20}{\micro\second} with a standard deviation of \qty{5}{\micro\second}. Data: \texttt{nw\_asia.scenarios()}.}\label{tab:nw:the-asia-pacific-map:designs}
\end{table}

\omcode{code/networks/12-the-asia-pacific-map/python/nw_asia.py}{70}{78}{Each design against the two parts of dissemination fairness: one tick's
spread, and whether the first member stays first.}\label{lst:nw:the-asia-pacific-map:designs}

\begin{dated}{2026-09}{The NSE co-location case}\label{dat:nw:the-asia-pacific-map:case}
Complaints to SEBI in 2015; SEBI's order of 30 April 2019 directing NSE to disgorge 624.89~crore rupees with interest at 12\% a year from
1~April 2014; the Securities Appellate Tribunal's later ruling setting the disgorgement aside, which SEBI appealed; SEBI's acceptance in July
2026 of NSE's proposals to settle the co-location and dark-fibre matters for about 1\,492~crore rupees (about 1\,224~crore for co-location);
and the Supreme Court's order of 18 September 2026 allowing the settlement and closing SEBI's appeals.
\end{dated}

The lesson for a trading firm is not only historical. Wherever a venue or a vendor distributes data over one connection per client (a
TCP feed, a websocket stream in chapter~20, a cloud relay in chapter~16), the order of the copies is a latency that no cable can fix, and the
questions of this section apply: in what order are the copies sent, does the order change, and who sits on the quieter server.

\section{Onshore China and its restrictions}

China's securities regulator has regulated co-location as a question of fairness. Its Administrative Rules for Program Trading in the
Securities Market, issued in May 2024 and in force from 8 October 2024, require an exchange that offers co-location to follow the principles of
``safety, fairness and reasonableness'' and to allocate the resources fairly, and brokers to keep trading fair between their clients; the
exchanges' accompanying notice singles out for monitoring accounts that reach 300 orders or cancellations a second, or 20\,000 in a day. A
consultation of April 2023 on trading-server co-location proposed that exchanges manage co-location in dedicated areas and keep the
communication latency from those areas to the front-end processors of the core trading system consistent: the equal-cable rule of
chapter~9, written as a latency requirement.

\begin{remark}[What the map leaves out]\label{rem:nw:the-asia-pacific-map:leftout}
The chapter does not place the mainland Chinese exchanges' buildings on the map: their co-location is reached through brokers, under the rules
above, and the site table would carry only a city. Chapters~22 to~26 add venues by asset class, with the same rule: no row without a
source, and a precision note when the source gives only a city.
\end{remark}

\begin{method}[Auditing a venue's dissemination design]\label{met:nw:the-asia-pacific-map:audit}
\begin{enumerate}
\item Find out how the feed is delivered in co-location: multicast, unicast (TCP), or both, as SGX will offer on its new platform.
\item If unicast, ask in what order copies are sent (by login, by connection, randomised per message), how many recipients share a sender,
  and whether some servers carry fewer recipients.
\item Measure: timestamp the same messages on several connections and servers (chapter~5), compute the spread and the rank correlation between
  messages.
\item Compare with the venue's published latency reports and fairness commitments; report a persistent order to the venue.
\end{enumerate}
\end{method}

\section{Tutorial: dissemination order}\label{tut:nw:the-asia-pacific-map:tut}

\begin{tutorial}
\textbf{Goal.} Extend the site table to Asia, then simulate unicast and multicast dissemination and measure their fairness.
\textbf{End state:} \cref{fig:nw:the-asia-pacific-map:map,fig:nw:the-asia-pacific-map:rank,fig:nw:the-asia-pacific-map:sorted} and
\cref{tab:nw:the-asia-pacific-map:pairs,tab:nw:the-asia-pacific-map:designs}.
\begin{tutsteps}
\item \textbf{Sites.} Five rows in \texttt{data/networks/sites.csv}, each with its precision in its name; \texttt{nw\_asia.pair\_rows()} computes
  \cref{tab:nw:the-asia-pacific-map:pairs}.
\item \textbf{Topology.} \texttt{nw\_asia.topology(m)} builds four servers of three ports of $m$ members and a secondary server of three ports of
  two; the costs are in \texttt{PARAMS}.
\item \textbf{One tick.} \texttt{firm\_dissem.unicast\_us} (\cref{lst:nw:the-asia-pacific-map:unicast}) and \texttt{multicast\_us} give every
  member's receive time; \texttt{fairness} summarises them.
\item \textbf{Many ticks.} \texttt{rank\_stability} and \texttt{race\_win} measure persistence and its value; \texttt{scenarios()}
  (\cref{lst:nw:the-asia-pacific-map:designs}) fills \cref{tab:nw:the-asia-pacific-map:designs}.
\end{tutsteps}
\textbf{What to change next.} Give the secondary server the same load as the others; halve the per-copy cost; put 20 members on each port and
check \cref{prop:nw:the-asia-pacific-map:seq}.
\end{tutorial}

\begin{omfigure}
\begin{tikzpicture}
\begin{axis}[width=11cm,height=6cm,xmin=0,xmax=100,ymin=0,ymax=175,xlabel={members, sorted by receive time (\%)},
  ylabel={receive delay (\si{\micro\second})},tick label style={font=\footnotesize},label style={font=\small},grid=major,
  grid style={gray!25},legend pos=north west,legend style={font=\footnotesize},table/col sep=comma]
\addplot[omDef,thick,no markers] table[x=pct,y=us]{figdata/networks/12-the-asia-pacific-map/sorted_unicast.csv};
\addplot[omMeth,thick,dashed,no markers] table[x=pct,y=us]{figdata/networks/12-the-asia-pacific-map/sorted_randomised.csv};
\addplot[omThm,thick,no markers] table[x=pct,y=us]{figdata/networks/12-the-asia-pacific-map/sorted_multicast.csv};
\legend{unicast,{unicast, randomised, no secondary},multicast}
\end{axis}
\end{tikzpicture}

\omcaption{Simulation: one tick's receive delays, members sorted from first to last, under three designs of
\cref{tab:nw:the-asia-pacific-map:designs}. Randomising changes who is where on the curve, not the curve; multicast flattens it. Data:
\texttt{fig\_asia.py}.}\label{fig:nw:the-asia-pacific-map:sorted}
\end{omfigure}

\section{Build: the dissemination model}\label{bld:nw:the-asia-pacific-map:dissem}

\begin{build}
\bfield{Purpose} A model of a venue's market-data distribution, to audit a documented design (chapter~22) and to explain a measured spread
between connections.
\bfield{Interface} \texttt{firm\_dissem}: \texttt{Server(name, ports, speed)}, \texttt{Member(server, port, rank)}, \texttt{members},
\texttt{unicast\_us(servers, c\_centre, c\_send, jitter, rng, shuffle)}, \texttt{multicast\_us}, \texttt{fairness}, \texttt{rank\_stability},
\texttt{race\_win}.
\bfield{Rules} Every cost is a parameter and every output is labelled a simulation; the model reproduces the documented order of sending, not a
venue's measured latency.
\bfield{Acceptance tests} \texttt{code/firm/dissem/tests/}: receive times by hand for a small topology, multicast flat, a fixed order stable
and a randomised one not, and the race probabilities at 0 and 50 microseconds.
\bfield{Stretch} TCP's own effects (a slow receiver's window stalling its port, retransmissions); a switch that replicates multicast port by
port.
\end{build}

\begin{omsources}
\item SEBI, Order in the matter of NSE colocation, 30 April 2019; SEBI circular CIR/MRD/DP/07/2015 on co-location and proximity hosting, 13 May
2015; Business Standard, 18 September 2026, on the Supreme Court's order.
\item JPX, connectivity services; HKEX, hosting services; SGX, Iris-ST FAQs (July 2026); ASX, Australian Liquidity Centre; NSE, co-location
capacity release (2025).
\item CSRC, Administrative Rules for Program Trading in the Securities Market (Trial), 2024, and the 2023 consultation on trading-server
co-location (as summarised by JunHe).
\end{omsources}

\section{Exercises}

\begin{exercise}[$\star$]\label{exo:nw:the-asia-pacific-map:1}
What is the one-way \omterm{def:nw:the-north-american-map:floor}{latency floor} between Singapore and Sydney in vacuum and in fibre?
\end{exercise}

\begin{exercise}[$\star$]\label{exo:nw:the-asia-pacific-map:2}
Using \cref{prop:nw:the-asia-pacific-map:seq} with $c = \qty{10}{\micro\second}$, how much later than the first member does the tenth member on
the same port receive a tick?
\end{exercise}

\begin{exercise}[$\star$]\label{exo:nw:the-asia-pacific-map:3}
Why does a venue's unpublished building address matter little for Tokyo--Hong Kong, and a great deal inside a data centre?
\end{exercise}

\begin{exercise}[$\star\star$]\label{exo:nw:the-asia-pacific-map:4}
A randomiser redraws the order of copies for every tick. Why does it not reduce the spread, and what does it change?
\end{exercise}

\begin{exercise}[$\star\star$]\label{exo:nw:the-asia-pacific-map:5}
In the model, why does the first member of a randomised tick usually sit on the secondary server?
\end{exercise}

\begin{exercise}[$\star\star$]\label{exo:nw:the-asia-pacific-map:6}
SGX will offer ITCH over both TCP and multicast in co-location. Which would you use for a latency-sensitive strategy, and what would you keep the
other for?
\end{exercise}

\begin{exercise}[$\star\star\star$]\label{exo:nw:the-asia-pacific-map:7}
\emph{Coding.} Compute the first member's advantage over the median member for 2, 5, 10 and 20 members per port, and compare with
\cref{prop:nw:the-asia-pacific-map:seq}.
\end{exercise}

\begin{exercise}[$\star\star\star$]\label{exo:nw:the-asia-pacific-map:8}
\emph{Find the flaw.} ``All co-located members had cables of the same length, so all received the tick-by-tick feed at the same time.''
\end{exercise}

\section{Problem: First to Log In}

\begin{problem}[{Weekend problem --- the value of logging in first}]\label{pb:nw:the-asia-pacific-map:1}
A venue sends its tick-by-tick feed over TCP from four distribution servers, each with three ports of $m$ members, and a lightly loaded
secondary server; the centre spends \qty{5}{\micro\second} per copy to a server, a port \qty{10}{\micro\second} per copy to a member. Members react
to a tick in \qty{20}{\micro\second} on average (standard deviation \qty{5}{\micro\second}).

\medskip\noindent\textbf{Part I --- One port.}
\begin{enumerate}
\item With $m = 10$, when do the first and the last member of the first server's first port receive a tick?
\item What does the median member of that port wait?
\item How does that change on the fourth server in the centre's order?
\item Which member of the whole topology receives first, and which last?
\end{enumerate}

\medskip\noindent\textbf{Part II --- The advantage.}
\begin{enumerate}[resume]
\item From the model, what is the first member's advantage over the median member for $m = 2, 5, 10, 20$?
\item Compare with \cref{prop:nw:the-asia-pacific-map:seq}.
\item How often does the first member beat the median member to the engine, tick after tick?
\item How large is the advantage against a race decided by a few microseconds (One Quant Book~11)?
\end{enumerate}

\medskip\noindent\textbf{Part III --- The remedies.}
\begin{enumerate}[resume]
\item What does a randomiser do to the spread, the rank correlation and the persistent win rate?
\item What does closing the secondary server to normal traffic add?
\item What does multicast do?
\item Which remedy would you ask the venue for, and why?
\end{enumerate}

\medskip\noindent\textbf{Part IV --- The verdict.}
\begin{enumerate}[resume]
\item State the \emph{named result}: the first member's advantage over the median member with ten members per port, with and without multicast.
\item Which of the model's parameters are facts from SEBI's order, and which are assumptions?
\item How would you measure the real spread on a venue's feed?
\item Why does the 2015 circular ask for similar latency with respect to exchange-provided infrastructure, rather than equal cables?
\item How does China's 2023 consultation phrase the same rule?
\item Where else does one-copy-per-client distribution appear in this book?
\item What is the difference between a latency advantage from a shorter path and one from a sending order?
\item In one sentence: why did logging in first matter?
\end{enumerate}
\end{problem}

\section{Interview questions}

\begin{interviewq}[$\star$ \iqroles{developer}]\label{iq:nw:the-asia-pacific-map:1}
What is the difference between unicast and multicast market data, and why do exchanges prefer multicast in co-location?
\end{interviewq}

\begin{interviewq}[$\star$ \iqroles{developer, trader}]\label{iq:nw:the-asia-pacific-map:2}
Roughly how far apart are Tokyo, Hong Kong and Singapore in light-time, and what does that mean for cross-market strategies?
\end{interviewq}

\begin{interviewq}[$\star\star$ \iqroles{developer}]\label{iq:nw:the-asia-pacific-map:3}
You receive the same TCP feed on two connections and one is consistently 40 microseconds later. How do you find out why?
\end{interviewq}

\begin{interviewq}[$\star\star$ \iqroles{trader, researcher}]\label{iq:nw:the-asia-pacific-map:4}
What does ``fair access'' to market data mean, and how would you test a venue against it?
\end{interviewq}

\begin{interviewq}[$\star\star$ \iqroles{developer}]\label{iq:nw:the-asia-pacific-map:5}
A vendor streams quotes to clients over websockets from one server. Who receives each update first?
\end{interviewq}

\begin{interviewq}[$\star\star\star$ \iqroles{developer, researcher}]\label{iq:nw:the-asia-pacific-map:6}
Design a market-data distribution for a new venue that cannot favour anyone. What do you measure to prove it?
\end{interviewq}
